\documentclass[10pt,a4paper]{amsart}
\usepackage[margin=19mm]{geometry}
\usepackage{amsmath,amssymb,mathtools}
\usepackage[scr=rsfs,cal=euler]{mathalfa}
\usepackage{xfrac}
\usepackage[unicode,hidelinks,bookmarks=false]{hyperref}
\newcommand{\ie}{i.e.\ }
\newcommand{\cf}{cf.\ }
\newcommand{\resp}{respectively\ }
\newcommand{\Subsection}{\subsection}
\providecommand{\nobreakdash}{\nobreak\hbox{-}}
\setlength{\emergencystretch}{2em}
\multlinegap0pt
\usepackage{xcolor}
\usepackage{amssymb}
\usepackage{tcolorbox}
\newtheorem{corollary}{Corollary}
\title{DT-PBO: an Interpretable Tree-based Surrogate Model for Preferential Bayesian Optimization}
\author{Nick Leenders, Thomas Quadt, Boris Cule, Roy Lindelauf, Herman Monsuur, Joost van Oijen, Mark Voskuijl}
\date{Source: arXiv:2512.14263v2 (CC BY 4.0)}
\begin{document}
\maketitle

\noindent\emph{Source and licence.} DT-PBO: an Interpretable Tree-based Surrogate Model for Preferential Bayesian Optimization. Version arXiv:2512.14263v2 (CC BY 4.0), distributed by arXiv under the Creative Commons Attribution 4.0 International licence, http://creativecommons.org/licenses/by/4.0/, as recorded in the arXiv licence metadata for this exact version and shown on the abstract page https://arxiv.org/abs/2512.14263v2. Full licence notice: \texttt{licenses/arxiv-2512.14263-CC-BY-4.0.txt}.\par\smallskip

\noindent\textit{Editorial scope.} Counted unit: the article's corollary cor:ancestor\_separation (source0.tex lines 774--778). The article's complete original proof of it is delivered with it (source0.tex lines 780--805, one \texttt{proof} environment of 26 lines). No mathematics is written, repaired or re-derived: the statement and the proof are the article's own lines, in the article's own notation, and no retained line is substituted or re-flowed.  No further source-local context is needed: every cross-reference of every retained line resolves inside the two delivered spans.  Nothing was omitted from any retained span, so the package carries no omission record.  No retained line contains a citation, so the wrapper carries no bibliography and no bibliography entry.  No retained line contains a figure environment, an image-inclusion command or drawing code, so no image is delivered and the package declares no asset dependency.  Editorial formatting only, in addition: an AMS \texttt{amsart} preamble that restates the article's own declarations and macros used by the retained lines, namely the environments \texttt{corollary} on separate counters, together with the standard packages those lines need.  The retained environments print in this document as Corollary 1 in order of appearance, whereas the article prints the same items with the numbering of its own full document.  The complete native source of the article as distributed in the arXiv e-print for this exact version is bundled unchanged as \texttt{sources/191087/source0.tex}.
\begin{corollary}[Exclusion of straddlers under ancestor separation]
\label{cor:ancestor_separation}
Let $x\in A_R$ and $x'\in A_L$. Under the defined tree split $s_A$ and utility score $f(x)$, it holds that $f(x)-f(x')>0$. 
Hence, every straddler is ordered solely by the parent split.
\end{corollary}

\begin{proof}
Recall the following definitions:
\begin{itemize}
\item $f(x)$ is the utility score of an item.
\item $s_A(x)$ represents the tree split itself. It assigns $-1$ if the item falls into the left child node ($A_L$) and $+1$ if it falls into the right child node ($A_R$). $\gamma_A$ is the "bonus" or "penalty" an item gets just for landing on the right or left side of the split. The total gap created by this split between the two sides is exactly $2\gamma_A$.
\item $h_A(x)$ is the residual or leftover utility. It represents the minor, local variations in an item's score that are not explained by the main split. 
\end{itemize}

Using these definitions, we can expand the difference in utility scores:
\[
f(x)-f(x')
=
\gamma_A s_A(x)+h_A(x)-\gamma_A s_A(x')-h_A(x')
=
2\gamma_A + h_A(x)-h_A(x').
\]
Since the residual utility difference is bounded by the total gap:
\[
|h_A(x)-h_A(x')|
\le
\sup_{u\in A_R,\;v\in A_L}|h_A(u)-h_A(v)|
<
2\gamma_A,
\]
we obtain $f(x)-f(x')>0$. Therefore, such pairs provide information about the already determined parent-level ordering only, and not about the choice among descendant splits within $A_L$ or $A_R$.
\end{proof}

\end{document}
